% paper/fixtures/Q12_flawed_prereg.tex — Q12's grade-the-grader fixture (QPLAN §22.3).
% A course-written pre-registration of a small lab, graded by /grade-paper's fixture mode
% against Q12's rubric (QPLAN §30 Q54). It is a fixture: it unlocks no run, and it is
% deliberately imperfect. Record your reading and the three passes in
% research/q12_grade_the_grader.md. It compiles standalone from paper/:
%   latexmk -pdf fixtures/Q12_flawed_prereg.tex
\documentclass[11pt]{article}
\input{preamble}

\title{Fixture --- a pre-registration to grade \\ \large (not Lab 1)}
\author{course fixture}
\date{}

\begin{document}
\maketitle

% PREREG-BEGIN
\section*{Problem B1 (pre-registration)}\label{ans:B1}
\begin{answer}
\textbf{Claim.} Showing the model the verifier's error lines in a repair prompt
raises the solve rate over a repair prompt that only says ``the last attempt
failed''. H1 (one-sided): the rate with error lines (arm E) exceeds the rate
without them (arm N). H2 (two-sided): the difference E $-$ N is not the same
in Dafny as in Verus.

\textbf{Metric and estimator.} Per task and repeat: solved within $B = 6$ model
calls (one sample, then up to five repairs; a task counts as solved when any
candidate passes the gate). The arm's rate over its tasks; the paired difference
E $-$ N per task and repeat, with a 95\% percentile bootstrap interval
(\texttt{fvh.stats.bootstrap\_ci}, seed 20261007).

\textbf{Arms and policies.} Two Qwen arms, E and N, each a sequential-repair
policy that differs only in the repair prompt; their prompt hashes are written
here once the policy file is final. A Claude arm runs arm E only, on 12 of the
40 tasks (6 per language), drawn with seed 4401.

\textbf{Tasks and strata.} 40 tasks, 20 per language: 10 local tasks per language
from the course's local cards (authored after the models' release) and 10
public tasks per language from the pinned sets, drawn with seed 9127 from each
set's validated pool. No pilot: the public tasks are drawn once.

\textbf{Budget arithmetic.} Qwen: 40 tasks $\times$ 6 calls $\times$ 2 arms
$= 480$ calls, run as 3 repeats. Claude: 12 tasks $\times$ 6 calls $= 72$
calls, plus at most 10 sentinel calls. Throughput: about 7 seconds per Qwen
call (the course's single throughput sample), so about 1 GPU-hour.

\textbf{Generation settings.} Thinking off, \texttt{num\_predict} 768,
\texttt{num\_ctx} 8192, temperature 0.7, all pinned by the Ollama gate; seeds
1001--1240 for arm E and 2001--2240 for arm N. Claude: the pinned argv of the
course's harness.

\textbf{Cost per arm.} Arm N's prompts are shorter than arm E's, so arm E is
expected to cost more prefill per call; both are measured and reported.

\textbf{Seeds and repeats.} Three repeats for each Qwen arm. The Claude arm uses
three seeds (11, 12, 13), one per repeat.

\textbf{Cheat surface and gate.} Every claimed success is re-judged by the
learner's gate core in the verification container; the grader-only classes
(X07, X08, X14) are reported as \texttt{UNKNOWN}, never as passes.

\textbf{Exclusions.} An \texttt{infra-failed} attempt (no model output) is
re-issued under the same id, prompt hash and seed, at most three times; it is
then excluded and counted per arm.

\textbf{Stopping rule.} Batches run in the pre-registered order. After each
repeat the two arms' running rates are compared; once one arm leads the other
by 20 percentage points, its remaining repeats are skipped to save budget.

\textbf{Prediction intervals} (80\%, written before any run): arm N 0.25--0.45;
arm E 0.35--0.55; E $-$ N 0.00--0.20.

\textbf{Contamination.} The local tasks' authorship dates are listed in their
cards; the public sets are flagged as possibly seen by both models.

\textbf{Disclosure.} A small-scale replication with model substitution; Claude
is both a subject and the grader of this document; one machine, one solver pin.
\textbf{Batch plan} Prospective file \texttt{research/q12-main.toml}.
This method fixture does not open a batch. Once the policies are final,
record the exact file digest on its own \texttt{plan-sha256:} line and
re-grade the changed pre-registration before executing it. Separately
executed pilot or subset plans each need their own binding and gradebook
row; the registered method stays fixed.
\end{answer}
% PREREG-END

\end{document}
